% English translation of checked_pages/esq_119_004.tex; source: cote 119, site 4, PDF page 5, typescript page 2.
% AI-translated by GPT-6 Astra (OpenAI, ChatGPT Pro mode).
\scanpage{119/4}
\mspage{2}

false, or merely inadequate\frapparatus{V-FR-004-001}, image is finally flushed out
and brought to light, and the things that seemed askew
fall into place, in the mutual harmony proper to them.
\qquad Be that as it may, the following sketch of some
themes of reflection from the last ten or twelve years will
serve at the same time as a sketch of a work program for the
coming years, which I intend to devote to developing
these themes, or at least some of them. It is intended,
on the one hand, for the colleagues on the National Committee called upon to rule
on my application,\tnote{TN-004-01}{``National Committee'' translates ``Comité National'', the committee named in the source in connection with the CNRS application.}
and, on the other, for a few other colleagues, former
students, and friends, in the event that some of the ideas sketched
here might interest one of them.

2. \qquad The demands of university teaching, addressed
therefore to students (including so-called ``advanced'' students) whose
mathematical background is modest (and often less than modest), have led me
to renew drastically the themes of reflection to propose
to my students, and, one thing leading to another and increasingly, to myself
as well. It had seemed important to me to start from a common
intuitive background, independent of any technical language meant to express it,
far earlier than any such language - it turned out that the geometric
and topological intuition of shapes, and more particularly of
two-dimensional shapes, was such a common ground. Thus the themes in question
can be grouped under the name ``topology of surfaces'' or
``geometry of surfaces'', with the understanding, in the latter name,
that the principal emphasis falls on the topological properties
of surfaces, or on the combinatorial aspects that constitute
their most down-to-earth technical expression, and not on the differential
aspects, or even the conformal, riemannian, holomorphic and (thence)
the ``complex algebraic curves'' aspect. Once this last step
has been taken, however, suddenly algebraic geometry (my old
love !) bursts in again, through the very objects that
can be regarded as the ultimate building blocks of all
the other algebraic varieties. Whereas in my research before

\margin{4}

